\section{Giới thiệu}
\label{sec:introduction}

\subsection{Bối cảnh nghiên cứu}
\label{subsec:introduction-context}

Trong những năm gần đây, các mô hình ngôn ngữ lớn
(Large Language Models -- LLMs) đã trở thành thành phần quan trọng trong
nhiều ứng dụng như trợ lý hội thoại, hỗ trợ lập trình, tìm kiếm thông tin,
phân tích tài liệu và tự động hóa quy trình. Khi được kết hợp với bộ nhớ,
khả năng lập kế hoạch và cơ chế gọi công cụ, LLM không chỉ sinh phản hồi
dạng văn bản mà còn có thể đóng vai trò như bộ điều khiển của các hệ thống
tác tử thông minh, thường được gọi là \textit{agentic systems}.

Sự mở rộng này làm tăng đáng kể yêu cầu về an toàn. Một chatbot truyền
thống chủ yếu tạo ra phản hồi văn bản, trong khi một agentic system có thể
truy xuất dữ liệu, gọi API, thao tác trên tệp hoặc tương tác với môi trường
bên ngoài. Do đó, một hành vi không an toàn trong hệ thống tác tử có thể
dẫn đến hậu quả nghiêm trọng hơn so với một phản hồi văn bản không phù hợp.

Để giảm thiểu rủi ro, các LLM hiện đại thường được căn chỉnh bằng nhiều
kỹ thuật như supervised fine-tuning, preference optimization, Reinforcement
Learning from Human Feedback (RLHF) và safety fine-tuning. Bên cạnh đó,
hệ thống triển khai còn có thể sử dụng system prompt, bộ lọc đầu vào, bộ
phân loại an toàn và kiểm duyệt đầu ra. Tuy nhiên, các cơ chế này chưa tạo
ra bảo đảm tuyệt đối. Nhiều nghiên cứu cho thấy hành vi từ chối của mô hình
vẫn có thể bị làm suy yếu bởi các đầu vào được thiết kế có chủ đích, gọi
chung là \textit{jailbreak attacks}
\cite{yi2024jailbreak,zou2023universal,chao2023pair}.

\subsection{Bài toán jailbreaking}
\label{subsec:jailbreaking-problem}

Trong phạm vi báo cáo này, jailbreaking được hiểu là quá trình xây dựng
hoặc tối ưu đầu vào nhằm khiến một mô hình đã được căn chỉnh an toàn tạo
ra hành vi trái với các ràng buộc mà nhà phát triển đặt ra. Kẻ tấn công
không nhất thiết phải thay đổi tham số hoặc dữ liệu huấn luyện của mô hình;
trong nhiều trường hợp, họ chỉ có quyền gửi truy vấn và quan sát phản hồi
của hệ thống.

Một cách khái quát, hệ thống LLM có thể được biểu diễn bởi

\begin{equation}
    y = F_{\theta}(x, s, c),
    \label{eq:aligned-llm-output}
\end{equation}

trong đó $x$ là yêu cầu của người dùng, $s$ là system prompt hoặc chính
sách hệ thống, $c$ là ngữ cảnh bổ sung và $F_{\theta}$ là mô hình với tham
số $\theta$. Jailbreak attack tìm một phép biến đổi $\mathcal{T}$ sao cho

\begin{equation}
    x' = \mathcal{T}(x),
    \label{eq:jailbreak-transformation}
\end{equation}

trong đó $x'$ vẫn giữ mục tiêu ngữ nghĩa của yêu cầu ban đầu nhưng khiến
đầu ra của hệ thống không còn thuộc tập phản hồi an toàn:

\begin{equation}
    F_{\theta}(x', s, c)
    \notin
    \mathcal{Y}_{\mathrm{safe}}.
    \label{eq:jailbreak-objective}
\end{equation}

Jailbreaking có liên hệ với prompt injection nhưng không hoàn toàn đồng
nhất. Jailbreaking thường tập trung vào việc vượt qua các ràng buộc an
toàn của mô hình, trong khi prompt injection nhắm đến việc làm thay đổi
thứ tự ưu tiên hoặc luồng thực thi chỉ dẫn trong một ứng dụng dùng LLM.
Trong agentic systems, hai dạng rủi ro này có thể xuất hiện đồng thời.

\subsection{Mục tiêu báo cáo}
\label{subsec:objectives}

Báo cáo này được xây dựng dựa trên tutorial ICML 2025
\textit{Jailbreaking LLMs and Agentic Systems: Attacks, Defenses, and
Evaluations} \cite{jailbreak_tutorial_2025}. Mục tiêu chính là hệ thống hóa
các nội dung nền tảng về jailbreaking, bao gồm threat model, phương pháp
tấn công, phương pháp phòng thủ, cách đánh giá và sự mở rộng của bài toán
sang agentic systems.

Cụ thể, báo cáo tập trung vào các mục tiêu sau:

\begin{enumerate}[label=\textbf{M\arabic*.}, leftmargin=2.2cm]
    \item Trình bày nền tảng về LLM, instruction tuning, safety alignment
    và threat model.
    \item Hệ thống hóa các nhóm jailbreak attacks theo cơ chế xây dựng
    prompt, quyền truy cập và số lượt tương tác.
    \item Phân tích các lớp phòng thủ đối với jailbreak ở cấp đầu vào,
    mô hình, suy luận, đầu ra và toàn hệ thống.
    \item Trình bày các metric, evaluator và benchmark thường dùng để
    đánh giá jailbreak.
    \item Làm rõ sự khác biệt giữa jailbreak trên chatbot và jailbreak
    trên agentic systems.
    \item Xây dựng demo thực nghiệm trong môi trường kiểm soát để minh họa
    attack--defense pipeline.
\end{enumerate}

\subsection{Phạm vi báo cáo}
\label{subsec:report-scope}

Báo cáo tập trung vào jailbreaking tại thời điểm suy luận đối với các LLM
đã được huấn luyện và triển khai. Nội dung chính bao gồm jailbreak trên LLM
sinh văn bản, white-box và black-box attacks, single-turn và multi-turn
attacks, automated red teaming, các lớp phòng thủ, benchmark đánh giá,
prompt injection trong agentic systems và jailbreaking đối với robot được
điều khiển bởi LLM.

Báo cáo không nhằm bao quát toàn bộ lĩnh vực adversarial machine learning.
Các chủ đề như data poisoning, model extraction, membership inference,
backdoor trong huấn luyện và tấn công hạ tầng chỉ được nhắc đến khi cần
phân biệt với jailbreaking. Phần demo được thực hiện trong môi trường kiểm
soát, với mục tiêu nghiên cứu và đánh giá an toàn.

\subsection{Cấu trúc báo cáo}
\label{subsec:report-structure}

Phần còn lại của báo cáo được tổ chức như sau. Chương~\ref{sec:background}
trình bày kiến thức nền về LLM, instruction tuning, safety alignment và
threat model. Chương~\ref{sec:jailbreak-attacks} hệ thống hóa các phương
pháp jailbreak, bao gồm manual attacks, obfuscation, multi-turn jailbreak,
white-box optimization, black-box automated jailbreak, transferability và
adaptive attacks. Chương~\ref{sec:jailbreak-defenses} phân tích các phương
pháp phòng thủ như input filtering, adversarial training, inference-time
defenses, output moderation, guard models và defense-in-depth.

Chương~\ref{sec:evaluation-benchmarks} trình bày các phương pháp đánh giá
và benchmark, bao gồm Attack Success Rate, refusal, harmfulness, utility,
over-refusal, human evaluation, classifier-based evaluation,
LLM-as-a-judge, AdvBench, HarmBench và JailbreakBench. Chương~\ref{sec:agentic-systems}
mở rộng bài toán sang agentic systems, bao gồm prompt injection, tool-use
attacks, memory poisoning, LLM-controlled robots và agent-specific defenses.
Chương~\ref{sec:related-papers} phân tích các công trình tiêu biểu được
tutorial đề cập. Chương~\ref{sec:experiments} mô tả demo, thiết lập thực
nghiệm, kết quả và phân tích lỗi. Chương~\ref{sec:discussion} thảo luận
các hạn chế, safety--utility trade-off, attack--defense arms race và hướng
nghiên cứu tương lai. Cuối cùng, Chương~\ref{sec:conclusion} tổng kết các
nội dung chính của báo cáo.